% RECOVERED_RESULT from II REV09. See MAPA_PROCEDENCIA.json.
% Explicit line selection and mechanical changes; no new proof.
\subsection{Angular coordinates, normalization of the turn, and units of action}
\label{v:ant:sec:ii-angulos}

The quantity of action and the angular coordinate belong to distinct
spaces. The section \(\hbar_{\rm ret}\) occupies the action line
\(\mathcal L_S\); \(\eta_{\rm ret}\) is its coefficient in the
declared basis. Angular representation acts on the dimensionless pair
\((\alpha,\eta_{\rm ret})\), after its construction:
\begin{equation}
 \mathcal P(\alpha,\eta_{\rm ret})
 =\bigl([1000\alpha]_{360},[2(\eta_{\rm ret}+\alpha)]_{360}\bigr).
 \label{v:ant:eq:ii-par-angular}
\end{equation}
Nonadic completion determines the factor 1000; bidirectional
closure determines the factor two. The circular chart of 360
units makes the partitions
\(12\cdot30=4\cdot90=3\cdot120\) visible. On the positive branch of
the article, the representatives are
\[
 A_{\deg}=1000\alpha,\qquad
 C^*_{\deg}=2(\eta_{\rm ret}+\alpha).
\]

\begin{proposition}[Shift of the base-thousand representation]
\label{v:ant:prop:ii-desplazamiento-angular}
Let \(\alpha=\sum_{n\geq1}b_n1000^{-n}\), with
\(b_n\in\{0,\ldots,999\}\), be the compatible expansion already generated.
The lifted coordinate \(A_{\deg}=1000\alpha\) satisfies
\[
 A_{\deg}=b_1+\sum_{n\geq1}b_{n+1}1000^{-n}.
\]
The representation preserves the block sequence from the second block onward;
the initial block becomes the integer part.
\end{proposition}
\begin{proof}
The series converges absolutely. Multiplying each term by 1000
and separating the initial term gives the identity. The operation
acts on the constructed expansion, with its blocks and memory;
the subsequent circular representation takes its class modulo 360.
\end{proof}

\begin{definition}[Charts of the angular coordinate]
The degree chart expresses a complete turn in 360 units.
The radian chart expresses it as \(2\pi\). The coordinate changes
and their inverses are
\[
 x=\frac{\pi}{180}A_{\deg},\qquad
 y=\frac{\pi}{180}C^*_{\deg},\qquad
 A_{\deg}=\frac{180}{\pi}x,\qquad
 C^*_{\deg}=\frac{180}{\pi}y.
\]
On a path with preserved winding number, the lifted coordinates
are used; the circular quotient expresses its class modulo
the corresponding turn.
\end{definition}

The definition characterizes degrees and radians through their exact relation
to the turn. The record also retains orientation, turns, and
transport memory. Representation of its angular class preserves
phase, whereas the lift distinguishes paths with the same
final phase. The physical reading of action requires specifying which of
these coordinates enters the integral.

\begin{proposition}[Covariance of the action reader]
For \(h_a=2\pi\hbar_a\), the action of a lifted coordinate
satisfies
\begin{equation}
 \mathcal S_a(\theta)
 =\hbar_a\theta_{\rm rad}
 =h_a\frac{\theta_{\deg}}{360}.
 \label{v:ant:eq:ii-accion-cartas}
\end{equation}
A positive turn expresses \(h_a\).
\end{proposition}
\begin{proof}
Substituting \(\theta_{\rm rad}=\pi\theta_{\deg}/180\) gives
\(\hbar_a\pi\theta_{\deg}/180
=(2\pi\hbar_a)\theta_{\deg}/360\). In one turn
\(\theta_{\deg}=360\), yielding \(h_a\).
\end{proof}

The periodicity of a representation is declared separately. If the
transported representation satisfies \(U(2\pi)=-I\), its operator
restoration occurs at \(4\pi\). This property of the lift
coexists with the per-turn normalization of
\eqref{v:ant:eq:ii-accion-cartas}. Each preserves its own object: integrated
quantity, circular phase, and represented state.

\subsection{Covariance of the channels and the incidence functional}

The channels are written equivalently as
\begin{equation}
 q_\pm=\exp[-(x\pm y)]
 =\exp\!\left[-\frac{\pi}{180}(A_{\deg}\pm C^*_{\deg})\right].
\label{v:ant:eq:ii-canales}
\end{equation}
\begin{proposition}[Recovery of the even and odd coordinates]
\label{v:ant:prop:ii-inversion-canales-angulares}
In the lifted real chart, the positive channels uniquely
determine the two coordinates:
\[
 A_{\deg}=-\frac{90}{\pi}\log(q_+q_-),\qquad
 C^*_{\deg}=\frac{90}{\pi}\log\frac{q_-}{q_+}.
\]
The exchange \((q_+,q_-)\mapsto(q_-,q_+)\) preserves
\(A_{\deg}\) and reverses \(C^*_{\deg}\).
\end{proposition}
\begin{proof}
The sum and difference of the logarithms of
\eqref{v:ant:eq:ii-canales} are, respectively,
\(-\pi A_{\deg}/90\) and \(-\pi C^*_{\deg}/90\).
The real logarithm is single-valued on the positive channels.
The product is symmetric, and the quotient is inverted when they are exchanged.
\end{proof}
In the twelve-position record, the antipodal involution
\(r\mapsto r+6\pmod {12}\) forms six pairs. The uniform
distribution of the odd coordinate over these pairs is
\(C^*_{\deg}/6\): summing their six contributions recovers
\(C^*_{\deg}\). This sectorial coordinate is distinguished from
the Planck constant \(h\), which belongs to the action line.

The conversion is applied only once to each argument. The quantity
\(\alpha\) remains dimensionless, and the quantity \(\hbar\)
continues to belong to the action line; their preceding coordinates
enter the map \(\mathcal P\).

In addition to the channels, the Barbero--Immirzi functional contains a
component bilinear in the angular representatives. Its transport is
\begin{equation}
 \frac{\pi}{180}A_{\deg}C^*_{\deg}
 =\frac{180}{\pi}xy.
 \label{v:ant:eq:ii-barbero-cartas}
\end{equation}
The identity follows by substituting both degree representatives
with their expressions in radians. The complete formula in analytic
coordinates is therefore
\[
 \gamma_{\HMT}
 =\frac{180}{\pi}xy
 +12\bigl[S_{90}(e^{-x+y})-S_{90}(e^{-x-y})\bigr]
 -\bigl[S_{120}(e^{-x+y})-S_{120}(e^{-x-y})\bigr].
\]
Equivalence of the two charts transports the bilinear component
and the exponential functions simultaneously. This transformation
makes explicit the precise role of the degree in the construction and allows
the evaluation to be reproduced without ambiguity in angular units.
